\honest{The Magnitude of His Own Folly}{Eliezer Yudkowsky}{2008}

Before I met the creationist who ran a funded AGI project, I used to argue with AGI hopefuls one at a time. I half-persuaded one of them, who had been much taken with evolutionary algorithms, that Friendly AI mattered and that finding the right fitness metric would not do. He lamented: ``Through my carelessness, I almost destroyed the world!'' That is a trap: ``What a villain I once was!''

I knew better, because in late 2002 I had looked back at my own AI proposals of 1997 and seen what they would have done, ``insofar as they were coherent enough'' to say. Everything fell into place at once; the doubts that had piled up behind a dam broke through together. I did not wonder how I could have been so stupid; I already knew how. To cry ``I almost destroyed the world!'' would have been too proud, confirming my own importance just when my ego ought to take a punch in the stomach. I had been so much less than I needed to be. To say that the gun was not really loaded, since I had no code ready to run, would also have softened the punch, since I had ``proposed and intended to build the gun'', load it, put it to my head and pull the trigger. Nor did I make an emotional drama of it, which would have spent the punch on tears. \nb{The archived record bears this out. In January 1999 Yudkowsky wrote on the Extropians mailing list of an open-source AI effort ``which I intend to oversee at some point.''}

I had not been updating for six years, and I had to stop: ``Halt, melt, and catch fire.'' Say ``I'm not ready.'' Say ``I don't know how to do this yet.'' In AGI those words cost a great deal of status, because researchers and the public want code, or at least someone ready to write code once funded. Without that, what distinguishes you from six billion other people who do not know how to build AGI? The field has no concept of ``I am trying to get from an incomplete map of FAI to a complete map of FAI.'' Name anyone else who says at once that they intend to build AGI, cannot yet because they do not know X, and are trying to figure out X. So there is a huge reluctance to say stop, and I felt it too, having absorbed the same attitude. \nb{This comes from the author's experience of the field; the post gives no case.}

My deeper reason for not stopping, even after I took up Friendly AI, was that I did not believe at gut level that ``Nature was allowed to kill me.'' The proverb says teenagers think they are immortal. They do not believe it literally, but their own death is not quite real to them, so seat belts do not compel them. I always wore my seat belt; I knew I could die. But raised in technophilia, I thought the Future was indestructible: even if nanotechnology could wipe out humanity, the Singularity would be too smart to be corrupted. \nb{In December 1996 Yudkowsky wrote that ``nanotechnology will turn the Earth into gray goo,'' and listed ``whether the Powers will be ethical'' among questions whose answers were ``Almost certainly'', ``of course'', and ``consider the alternative.''}

``The thought you cannot think controls you more than thoughts you speak aloud.'' But we flinch only from fears that are real to us. AGI researchers can picture headlines saying their work has been upstaged, and those who have started companies know they can run out of venture capital. I guess that ``Oops'' followed by the thud of six billion bodies, by their own hands, is not real to them in the same way. They fear being scooped more than destroying the world by their own mistake, and so they say that if you delay for safety, others will beat you. I mark this as a guess (``It is unsafe to say what other people are thinking''), and I admit: ``I, too, used to say things like that.'' \nb{In \textsc{Creating Friendly AI} (2001) Yudkowsky wrote that a controlled ascent should be ``fast enough not to give a significant speed advantage to a non-Friendliness-aware AI project.''}

Then my technophilia broke. I was ``not being graded on a curve'': even if you follow the rules of science, are a nice person and run the best project of all, Nature can still kill you. My gaze shook free of rivals and I saw the sheer blank wall. My careful arguments for going on at full speed, and my reasons to tolerate the risk of a fundamental mistake, would not have stopped that risk from killing me. Nature is allowed to say ``So what?'' to the best argument. Only my ignorance of the rules let me argue for going ahead without knowing them, because if you do not know the rules you cannot model the penalty for ignorance. Others still say they do not know that their plan cannot work; when I explain how small the target is in the search space, they ask, ``How can you be so sure I won't win the lottery?'' \nb{\textsc{Creating Friendly AI} had argued that many harmless mistakes would show up before a catastrophic one, so that watchers who saw none ``can probably be fairly confident that none has ever occurred.''}

The only safe course in my ignorance was to say: ``I will not proceed until I know positively that the ground is safe.'' The clever arguments for stepping on ground that may hide a landmine sound much less clever once you look where you meant to step and see the bang. You can do everything you are supposed to do, and Nature can still kill you. ``That was when my last trust broke. And that was when my training as a rationalist began.''
